% !TeX spellcheck = en_US
\documentclass[french]{yLectureNote}

\title{Algèbre linéaire 3*}
\subtitle{Physique}
\author{Paulhenry Saux}
\date{\today}
\yLanguage{Français}
%lombardi@math.univ-toulouse.fr
\professor{E.Lombardie}
\usepackage{graphicx}%----pour mettre des images
\usepackage[utf8]{inputenc}%---encodage
\usepackage{geometry}%---pour modifier les tailles et mettre a4paper
%\usepackage{awesomebox}%---pour les boites d'exercices, de pbq et de croquis ---d\'esactiv\'e pour les TP de PC
\usepackage{tikz}%---pour deiffner + d\'ependance de chemfig
% \usepackage{tabularx}%---pour dimensionner automatiquement les tableaux avec variable X
\usepackage{awesomebox}%---Pour les boites info, danger et autres
\usepackage{menukeys}%---Pour deiffner les touches de Calculatrice
\usepackage{fancyhdr}%---pour les en-t\^ete personnalis\'ees
\usepackage{blindtext}%---pour les liens
\usepackage{hyperref}%---pour les liens (\`a mettre en dernier)
\usepackage{caption}%---pour la francisation de la l\'egende table vers Tableau
\usepackage{pifont}
\usepackage{array}%---pour les tableaux
\usepackage{lipsum}
\usepackage{yFlatTable}
\usepackage{multicol}
\newcommand{\Lim}[1]{\lim\limits_{\substack{#1}}\:}
\renewcommand{\vec}{\overrightarrow}
\newcommand{\N}[0]{\mathbb{N}}
\newcommand{\R}[0]{\mathbb{R}}
\newcommand{\dd}{\mathrm{d}}
\newcommand{\norm}[1]{||\vec{#1}||}
\newcommand{\fo}{\psi(\vec{r},t)}
\newcommand{\foe}{\psi(\vec{r},t)\*}
\newcommand{\HH}{\hat{H}}
\newcommand{\hb}{\hbar}
\newcommand{\lap}{\nabla^2}
\newcommand{\lapcc}{\frac{\partial^2 }{\partial x^2}+\frac{\partial^2 }{\partial y^2}+\frac{\partial^2 }{\partial z^2}}
\newcommand{\E}{\vec{E}(M)}
\newcommand{\B}{\vec{B}(M)}
\newcommand{\V}{V(M)}
\newcommand{\er}{\vec{e_{\rho}}}
\newcommand{\ep}{\vec{e_{\varphi}}}
\newcommand{\pip}{\pi^+}
\newcommand{\pim}{\pi^-}
\newcommand{\mv}{\mathcal{V}}
\DeclareMathOperator{\Div}{Div}
\newcommand{\rot}{\vec{\mathrm{rot}}}
\newcommand{\grad}{\vec{\mathrm{grad}}}
\begin{document}
\setcounter{chapter}{2}
\chapter{Espaces Hermitiens}
\section{Produit scalaire, norme}
Dans ce chapitre, on travaille sur des espaces vectoriels sur $\mathbb{C}$.

\begin{definition}
    Soit $E$ un espace vectoriel sur $\mathbb{C}$ de dimension finie ou non. Une application $C: E \times E \to \mathbb{C}$ est un produit scalaire hermitien si et seulement si :
    \begin{enumerate}
        \item $C(x+x', y) = C(x, y) + C(x', y)$ (additivité à gauche)
        \item $C(x, y+y') = C(x, y) + C(x, y')$ (additivité à droite)
        \item $C(x, \lambda y) = \lambda C(x, y)$ (linéarité à droite pour les scalaires)
        \item $C(\lambda x, y) = \bar{\lambda} C(x, y)$ (anti-linéarité à gauche pour les scalaires)
        \item $C(y, x) = \overline{C(x, y)}$ (symétrie hermitienne)
        \item $\forall x \in E, C(x, x) \ge 0$ (positivité) et $\left(C(x, x) = 0 \implies x = 0_E\right)$ (définie)
    \end{enumerate}
\end{definition}

\warningInfo{Attention aux conventions}{La définition ci-dessus est la convention (linéaire à droite, anti-linéaire ou demi-linéaire à gauche).}

On peut en déduire quelques conséquences directes 
\begin{itemize}
    \item L'application $y \mapsto C(x, y)$ est linéaire.
\item  L'application $x \mapsto C(x, y)$ est semi-linéaire (ou anti-linéaire) : les scalaires sortent conjugués.
\item  Pour tout $x \in E$, $C(x, x) \in \mathbb{R}$ en vertu de la symétrie hermitienne ($C(x,x) = \overline{C(x,x)}$). Cela donne tout son sens à la condition de positivité.
\end{itemize}:
\begin{definition}[Espace hermitien]
    Si $\dim E < +\infty$ et que $E$ est muni d'un produit scalaire hermitien, alors $E$ est un \textbf{espace hermitien}. Dans le cas contraire (dimension infinie), il est qualifié de \textbf{préhilbertien complexe}.
\end{definition}

Dans l'espace hermitien type $\mathbb{C}^n$ muni du produit scalaire canonique, on a :
$$(x|y) = \sum_{k=1}^n \bar{x}_k y_k \implies (x|x) = |x_1|^2 + |x_2|^2 + \dots + |x_n|^2$$

\begin{definition}[Notation]
    Désormais, on note $C(x, y) = (x|y) = \langle x|y \rangle$.
\end{definition}

\begin{definition}[Norme]
    La norme hermitienne associée est $\|x\| = \sqrt{(x|x)}$.
\end{definition}

\begin{theorem}[Inégalité de Cauchy-Schwarz]
    Pour tous $x, y \in E$ :
    $$|(x|y)| \leq \|x\| \|y\|$$
    Avec égalité si et seulement si la famille $(x, y)$ est liée.
\end{theorem}
\begin{theorem}[Propriétés de la norme]
\begin{itemize}
    \item $\|\lambda x\| = |\lambda| \|x\|$ (Homogénéité : attention au module !)
    \item $\|x+y\| \leq \|x\| + \|y\|$ (Inégalité triangulaire)
\end{itemize}
\end{theorem}

\begin{proposition}[Identité de polarisation]
    L'identité de polarisation en complexe permet de reconstituer le produit scalaire à partir de la norme :
    $$(x|y) = \frac{1}{4} \big(\|x+y\|^2 - \|x-y\|^2 - i\|x+iy\|^2 + i\|x-iy\|^2\big)$$
\end{proposition}

\subsection{Orthogonalité}

\begin{definition}
    Deux vecteurs $x$ et $y$ sont orthogonaux (noté $x \perp y$) si et seulement si $(x|y) = 0$.
\end{definition}

\warningInfo{Asymétrie visuelle}{Puisque $(x|y) = 0 \iff \overline{(y|x)} = 0 \iff (y|x) = 0$, la relation d'orthogonalité reste symétrique, même si le produit scalaire ne l'est pas.}

\begin{theorem}[Théorème de Pythagore]
    Si $x \perp y$, alors $\|x+y\|^2 = \|x\|^2 + \|y\|^2$.
\end{theorem}

\criticalInfo{Équivalence fausse en complexe}{En réel, Pythagore est une équivalence. En complexe, ce n'est qu'une implication. En effet, $\|x+y\|^2 = \|x\|^2 + \|y\|^2 + 2\text{Re}((x|y))$. Donc $\|x+y\|^2 = \|x\|^2 + \|y\|^2 \iff \text{Re}((x|y)) = 0$. On peut avoir une partie réelle nulle sans que le produit scalaire soit nul (ex: $(x|y) = i$).}

\begin{proposition}
    Soit $(e_1, \dots, e_n)$ une \textbf{base orthonormale} (BON) de $E$. Pour tout $x \in E$ :
    $$x = (e_1|x)e_1 + \dots + (e_n|x)e_n$$
    Les coordonnées de $x$ dans cette base sont $x_i = (e_i|x)$. 
\end{proposition}

\criticalInfo{Attention à l'ordre des facteurs}{Dans notre convention (linéaire à droite), le produit scalaire doit impérativement s'écrire $(e_i|x)$. Si on écrit $(x|e_i)$, on a $\bar{x}_i$ (le conjugué), ce qui fausse les calculs de coordonnées.}

% \warningInfo{Gram-Schmidt}{L'algorithme de Gram-Schmidt s'applique à l'identique pour normaliser et orthogonaliser une base. Il faut cependant conjuguer le premier terme du produit scalaire lors des projections : $u_k = v_k - \sum_{j=1}^{k-1} \frac{(u_j|v_k)}{\|u_j\|^2}u_j$.}

\begin{definition}[Orthogonal d'une partie]
    L'orthogonal d'une partie $A \subset E$, noté $A^\perp$, est défini par :
    $$A^\perp = \{x \in E \mid \forall a \in A, (x|a) = 0\}$$
\end{definition}

\begin{proposition}
    Si $F$ est un sous-espace vectoriel de $E$ de dimension finie, alors :
    $$E = F \oplus F^\perp, \quad \dim E = \dim F + \dim F^\perp \quad \text{et} \quad F^{\perp\perp} = F$$
\end{proposition}

\begin{proposition}
    Soit $(e_1, \dots, e_m)$ une BON d'un sous-espace $F$. La projection orthogonale de $x$ sur $F$ est :
    $$\pi_F(x) = \sum_{k=1}^m (e_k|x)e_k$$
\end{proposition}

\section{Endomorphisme adjoint}

\begin{definition}
    Soit $E$ un espace hermitien et $f \in \mathcal{L}(E)$. Il existe un unique endomorphisme de $E$, noté $f^*$, appelé \textbf{adjoint de $f$}, vérifiant :
    $$\forall x, y \in E, \quad (f(x)|y) = (x|f^*(y))$$
    Si $A = \text{Mat}_{\mathcal{B}}(f)$ dans une BON $\mathcal{B}$, alors la matrice de $f^*$ dans cette même base est la matrice adjointe :
    $$A^* = \overline{{}^tA}$$
\end{definition}

\begin{proposition}
    Propriétés de l'adjoint :
    \begin{itemize}
        \item $f^{**} = f$
        \item $(f+g)^* = f^* + g^*$
        \item $(\lambda f)^* = \bar{\lambda}f^*$ (L'application $f \mapsto f^*$ est \textbf{semi-linéaire})
        \item $\text{rg}(f) = \text{rg}(f^*)$
        \item $\det(f^*) = \overline{\det(f)}$
        \item $(f \circ g)^* = g^* \circ f^*$ (Attention au retournement de l'ordre)
    \end{itemize}
\end{proposition}

\section{Groupe unitaire}

\begin{proposition}
    Soit $E$ un espace hermitien et $f \in \mathcal{L}(E)$. Les assertions suivantes sont équivalentes :
    \begin{itemize}
        \item $\forall x \in E, \|f(x)\| = \|x\|$ (Isométrie vectorielle)
        \item $\forall x, y \in E, (f(x)|f(y)) = (x|y)$ (Conservation du produit scalaire)
        \item L'image de toute BON par $f$ est une BON.
        \item $A^*A = I_n \iff \overline{{}^tA}A = I_n$ (où $A$ est la matrice de $f$ dans une BON)
    \end{itemize}
\end{proposition}

\begin{definition}[Matrice unitaire]
    Une matrice carrée $A \in \mathcal{M}_n(\mathbb{C})$ vérifiant $A^*A = I_n$ est dite \textbf{unitaire}. L'ensemble de ces matrices forme le groupe unitaire, noté $\mathcal{U}_n(\mathbb{C})$.
\end{definition}

\begin{proposition}
    Si $A \in \mathcal{U}_n(\mathbb{C})$, alors :
    * $|\det(A)| = 1$. Le déterminant est un nombre complexe de module 1 (de la forme $e^{i\theta}$).
    * Si $\lambda \in \text{Sp}(A)$, alors $|\lambda| = 1$. Les valeurs propres sont toutes sur le cercle unité du plan complexe.
\end{proposition}

\begin{definition}
    Le sous-groupe de $\mathcal{U}_n(\mathbb{C})$ constitué des matrices de déterminant égal à 1 est appelé \textbf{groupe spécial unitaire}, noté $\mathcal{SU}_n(\mathbb{C})$.
\end{definition}

\tipsInfo{Lien avec le cas réel}{Le groupe orthogonal spécial $\mathcal{SO}_n(\mathbb{R})$ s'obtient simplement par l'intersection : $\mathcal{SO}_n(\mathbb{R}) = \mathcal{SU}_n(\mathbb{C}) \cap \mathcal{M}_n(\mathbb{R})$.}

% \subsubsection*{Structure des matrices de $\mathcal{SU}_2(\mathbb{C})$}
% Une matrice $A \in \mathcal{M}_2(\mathbb{C})$ appartient à $\mathcal{SU}_2(\mathbb{C})$ si et seulement s'il existe $a, b \in \mathbb{C}$ tels que $|a|^2 + |b|^2 = 1$ sous la forme :
% $$A = \begin{pmatrix} a & -\bar{b} \\ b & \bar{a} \end{pmatrix}$$
% En paramétrage polaire, cela donne l'écriture trigonométrique suivante (avec $\alpha, \theta, \phi \in \mathbb{R}$) :
% $$A = \begin{pmatrix} \cos(\alpha)e^{i\theta} & -\sin(\alpha)e^{-i\phi} \\ \sin(\alpha)e^{i\phi} & \cos(\alpha)e^{-i\theta} \end{pmatrix}$$

\section{Diagonalisation des endomorphismes auto-adjoints}

\begin{definition}
    Un endomorphisme $f$ est dit \textbf{auto-adjoint} (ou \textbf{hermitien}) si $f^* = f$.
    Cela se traduit par : $\forall x,y \in E, (f(x)|y) = (x|f(y))$.
    Dans une BON, sa matrice vérifie $A^* = A$ (matrice hermitienne).
\end{definition}

\begin{theorem}[Théorème spectral]
    Si $f$ est un endomorphisme auto-adjoint d'un espace hermitien $E$, alors :
    \begin{enumerate}
        \item Toutes les valeurs propres de $f$ sont \textbf{réelles} ($\text{Sp}(f) \subset \mathbb{R}$).
        \item Les sous-espaces propres de $f$ sont orthogonaux deux à deux.
        \item $f$ est diagonalisable dans une \textbf{base orthonormale} de vecteurs propres.
    \end{enumerate}
\end{theorem}
\end{document}